%----------------------------------------------------------------------
\subsection{Session 10: the block count is Harer--Zagier, the environment
is healthy on the short horizon, the exact size-biased law, and the
Markov-renewal form of the survival problem}\label{sec:s10}

Session 10 was to fix the exact statement behind
Problem~\ref{prob:survival} (TODO~10a, step~0) and choose a route.  The
first finding is negative and structural: \emph{the environment cannot
be avoided} --- without merges the parity of the flip count from a start
with $\min(a_0,b_0)=\varepsilon$ decays on the time scale
$\varepsilon^{-2}$, not $\varepsilon^{-1}$, and the budget of
Corollary~\ref{cor:kappafree} forces $\varepsilon\ll d^{-1/2}$, so a
proof must use the recharge of small marked coordinates by
environment blocks (\S\ref{sec:noenv}).  The second finding removes the
obstacle this creates.  The number of curves $B_t$ after $t$ chords has
an \emph{exact} law --- it is the Harer--Zagier vertex-count
distribution of a random gluing of a $2t$-gon --- with
$\E[3^{B_t}]=3+4t+2t^2$; a Chernoff bound then gives $B_t\le C\log d$
simultaneously for all $t\le9d$ except with probability $d^{-\gamma}$
(\S\ref{sec:blockcount}).  On that event every state on the horizon
has $\sum p_i^2\ge1/(C\log d)$ and an environment block of mass
$\ge(1-s)/(C\log d)$: the ``no-dust'' question is closed, up to
logarithms, in exactly the form the short horizon needs.  The same hook
computation gives the exact law of the mass of the curve through a
uniform point at every time (\S\ref{sec:sizebiased}; this contains
Theorem~\ref{thm:exactonepoint} and closes TODO~10b(b)).  Finally the
survival problem is put in Markov-renewal form
(\S\ref{sec:markovrenewal}): an exact operator identity for the
generating function $\sum_d\Phi_dz^d$ separates the two remaining
contents --- moments of the first-flip time (traps) and invertibility
of $I+K_z$ (cancellation) --- and a conditional decay proposition
states precisely what is left.

\subsubsection{The environment cannot be avoided}\label{sec:noenv}

Consider the marked pair alone, with merges suppressed (an adversarial
``dust'' environment): per step, $a\to a\sqrt U$ with probability
$a^2$, $b\to b\sqrt U$ with probability $b^2$, a flip
$(a,b)\to(\alpha+\beta,s-\alpha-\beta)$ with probability $2ab$, and
nothing otherwise.  Since merges only increase marked coordinates and
are the only moves that do, this is the chain that any argument using
the environment solely through upper bounds would have to control.
Its scaling is elementary: all rates are quadratic in the coordinates,
so from $(a_0,b_0)=(\varepsilon,\varepsilon)$ nothing happens for
$\asymp\varepsilon^{-2}$ steps, and from $(c_0,\varepsilon)$ the large
coordinate decays like $t^{-1/2}$ before a flip (rate $2\varepsilon b$)
can occur, the first flip arriving only at $t\asymp\varepsilon^{-2}$.
Measured (\texttt{s10\_nomerge}, $\Phi_t=\E[(-1)^{N_t}]$):
from $(\tfrac14,\tfrac14)$, $\Phi_t=0.130,\,0.031,\,0.0070,\,0.0024$
at $t=8,16,32,64$ (then noise floor $10^{-3}$; $t^2\Phi_t\approx8$--$10$
on $16\le t\le64$) --- a healthy $1/t^2$-like decay without any
environment; but from $(\tfrac14,0.02)$,
$\Phi_t=0.53,\,0.25,\,0.088,\,0.040$ at $t=64,256,1024,4096$, and from
$(0.02,0.02)$, $\Phi_t=0.90,\,0.68,\,0.30,\,0.13$ at
$t=64,256,1024,8192$: decay only on the scale $\varepsilon^2t$, and
slow even there.  In the budget of Corollary~\ref{cor:kappafree} the
trap-entry flux forces $\varepsilon\ll d^{-1/2}$, hence
$\varepsilon^2d\ll1$: a bound depending on the environment only
through upper bounds is worthless.  \emph{Status: COMPUTED (the scaling
statement is trivial).  Consequence: any proof of
Problem~\ref{prob:survival} must lower-bound the masses of environment
blocks over the horizon --- the recharge mechanism of
Remark~\ref{rem:renewal}.  The next subsection supplies exactly this.}

\subsubsection{The exact law of the number of curves}\label{sec:blockcount}

Let $B_t$ be the number of curves after $t$ iid chords (the total
block count of the marked chain from the one-block start; it does not
depend on the marks).

\begin{theorem}[the block count is Harer--Zagier]\label{thm:blockcount}
For every $x$,
\[
\sum_{t\ge0}\E\bigl[x^{B_t}\bigr]z^t
=\frac1{2z}\Bigl[\Bigl(\frac{1+z}{1-z}\Bigr)^{x}-1\Bigr],
\qquad\text{i.e.}\qquad
\Pr[B_t=m]=[z^{t+1}]\,\frac{(2\operatorname{artanh}z)^m}{2\,m!},
\]
so $B_t\equiv t+1\pmod2$, $\Pr[B_t=t+1-2g]=\varepsilon_g(t)/(2t-1)!!$
with $\varepsilon_g(t)$ the Harer--Zagier numbers \cite{HZ86}, and for
integer $x\ge1$ the expectation $\E[x^{B_t}]$ is a polynomial in $t$ of
degree $x-1$:
\[
\E[2^{B_t}]=2+2t,\qquad \E[3^{B_t}]=3+4t+2t^2,\qquad
\E[4^{B_t}]=4+\tfrac{20}3t+4t^2+\tfrac43t^3=\tfrac43(t+1)(t^2+2t+3),\qquad
\Pr[B_t=1]=\frac{1+(-1)^t}{2(t+1)}.
\]
\end{theorem}

\begin{proof}
The $2t$ endpoints are iid uniform on the circle, a.s.\ distinct; by
exchangeability, conditionally on their positions the pairing into
chords is uniform among the $(2t-1)!!$ pairings.  Refining $\gamma$ at
the endpoints does not change the curve partition
(Lemma~\ref{lem:splitmerge}), so $B_t$ is the number of cycles of
$\gamma_{2t}\tau$, with $\gamma_{2t}$ the rotation of the $2t$
endpoints and $\tau$ the pairing involution --- the number of vertices
of the surface obtained by gluing the sides of a $2t$-gon along a
uniform pairing.  Harer--Zagier's formula \cite{HZ86} is
$1+2\sum_{t\ge0}\bigl[\sum_g\varepsilon_g(t)N^{t+1-2g}\bigr]
z^{t+1}/(2t-1)!!=\bigl((1+z)/(1-z)\bigr)^N$; dividing by $(2t-1)!!$ turns
the bracket into $\E[N^{B_t}]$, which is the displayed identity for
$x=N$, and both sides are polynomials in $N$ for each $t$.  Expanding
$((1+z)/(1-z))^x=\exp(2x\operatorname{artanh}z)$ in powers of $x$ gives
the law; the polynomials are the coefficients of
$(1-z)^{-x}\sum_k\binom{x}{2k+1}z^{2k}$.
\end{proof}

\emph{Status: PROVED (the statement is classical once the chord
diagram is identified; the identification is the content).  A second
proof runs through the finite-$n$ hook machinery of
Theorem~\ref{thm:exactonepoint}: $x^{c(\sigma)}=\sum_\lambda
s_\lambda(1^x)\chi_\lambda(\sigma)$, so
$\E[x^{c(\sigma_t)}]=\sum_{k}(-1)^k r_k^{\,t}\,s_{(n-k,1^k)}(1^x)$ with
$s_{(n-k,1^k)}(1^x)=x(x{+}1)\cdots(x{+}n{-}k{-}1)\cdot(x{-}1)\cdots(x{-}k)
/\bigl(n\,(n{-}k{-}1)!\,k!\bigr)$ by hook-content; for integer $x$ only
$k\le x-1$ survive, e.g.\ $\E[2^{c(\sigma_t)}]=(n+1)-(n-1)r_1^{\,t}$,
$\E[3^{c(\sigma_t)}]=\tfrac12(n+1)(n+2)-(n^2-1)r_1^{\,t}+\tfrac12(n-1)(n-2)r_2^{\,t}$,
whose $n\to\infty$ limits are $2+2t$ and $3+4t+2t^2$
(\texttt{s10\_blocks.py}, exact rationals, Lagrange-extrapolated in
$1/(n-1)$, gives the polynomials for $x\le5$).  VERIFIED: the full law
against a direct simulation of the split--merge chain
(\texttt{s10\_blocksim.py}, $2\times10^5$ trials, $t\le6$): every
probability agrees to $\pm0.001$, e.g.\ $\Pr[B_6=1,3,5,7]=
0.1429,0.6222,0.2222,0.0127$ predicted vs.\ $0.1431,0.6230,0.2212,0.0127$
measured.}

\begin{corollary}[block count over the horizon]\label{cor:horizon}
For all $T\ge1$ and $K\ge1$,
\[
\Pr\Bigl[\max_{0\le t\le T}B_t\ge K\Bigr]\;\le\;3^{\,1-K}(T+1)^3 .
\]
In particular, for $\gamma>0$ and
$K_\gamma(T):=\lceil(3+\gamma)\log_3(T+1)\rceil+1$, the event
$H_T=\{B_t<K_\gamma(T)\ \forall t\le T\}$ has probability
$\ge1-(T+1)^{-\gamma}$, and on $H_T$, for every $t\le T$: the total number of
blocks is $<K_\gamma(T)$, $\sum_ip_i(t)^2\ge1/K_\gamma(T)$, the
environment has fewer than $K_\gamma(T)$ blocks and hence a block of
mass $\ge(1-s_t)/K_\gamma(T)$.
\end{corollary}

\begin{proof}
Markov with $x=3$: $\Pr[B_t\ge K]\le3^{-K}(3+4t+2t^2)\le3^{1-K}(t+1)^2$;
sum over $t\le T$.  With $K=K_\gamma(T)$,
$3^{1-K}\le(T+1)^{-3-\gamma}$.  The consequences are Cauchy--Schwarz
($\sum p_i^2\ge1/\#\text{blocks}$) and pigeonhole.
\end{proof}

\emph{Status: PROVED.  In Problem~\ref{prob:survival} the start
$\Gamma_0\sim\mu_j$ ($j\le8d$) followed by $d$ steps is the chain from
one block over at most $9d$ steps, so $H_{9d}$ has probability
$\ge1-(9d+1)^{-\gamma}$ for any fixed $\gamma$ at the price
$K=O_\gamma(\log d)$.  This is the ``no-dust'' estimate that
\S\ref{sec:upper} and the session-8/9 handoffs listed as the open
ingredient of every route: over the short horizon it holds for free,
up to a logarithm, because the chain simply has not had time to create
many blocks --- $\E[B_t]=[z^t]\,(1+z)\operatorname{artanh}(z)/(z(1-z))=\log t+O(1)$
and the pgf makes the tail geometric.  (At stationarity the count is infinite and the analogous
statement is false as stated; the short horizon of
Theorem~\ref{thm:factor} is again what makes the iid model tractable.)
Since the required exponent in Corollary~\ref{cor:kappafree} is only
$\alpha>2/9$, losing logarithms is harmless.}

\subsubsection{The exact law of the block of a uniform point}
\label{sec:sizebiased}

\begin{theorem}[size-biased block mass at all times]\label{thm:sizebiased}
Let $U$ be uniform on the circle, independent of the chords, and let
$\zeta_m$ be the mass of the curve through $U$ after $m$ chords (chain
from one block).  Then
\[
\Pr[\zeta_m\in dx]=\bigl(1-(1-2x)^m\bigr)\,dx\ \text{ on }(0,1),
\qquad
\Pr[\zeta_m=1]=\frac{1+(-1)^m}{2(m+1)} ,
\]
and consequently, for every $r\ge1$,
\[
\E\Bigl[\sum_ip_i(m)^r\Bigr]=\E[\zeta_m^{\,r-1}]
=\frac1r+\int_0^1(1-x^{r-1})(1-2x)^m\,dx .
\]
By rotation invariance the same law governs the mass of the curve
through any fixed point, in particular through $u$ (and through $v$).
At finite $n$: $\Pr[|C_{\sigma_m}(1)|=j]=(1-r_j^{\,m})/n$ for $j<n$ and
$\Pr[|C_{\sigma_m}(1)|=n]=\frac1n\sum_{k=0}^{n-1}r_k^{\,m}$, with
$r_k=1-2k/(n-1)$; equivalently
$(-1)^k\langle\chi_{(n-k,1^k)},\sum_ic_i^{\,r}\rangle=n^{r-1}-k^{r-1}$
for $1\le k\le n-1$ and $\sum_{j\le n}j^{r-1}$ for $k=0$.
\end{theorem}

\begin{proof}
For $F$ on $\{1,\dots,n\}$ put $G_F(\sigma)=\sum_{u}F(|C_\sigma(u)|)
=\sum_ic_iF(c_i)$.  The exponential formula with one marked cycle, as
in the proof of Theorem~\ref{thm:exactonepoint} (weight $jF(j)$ for a
marked cycle of length $j$ replaces $j\cdot j$ there), gives
\[
\sum_nx^n\sum_kt^k\langle\chi_{(n-k,1^k)},G_F\rangle
=\frac1{1+t}\Bigl[\sum_{j\ge1}F(j)\bigl(1-(-t)^j\bigr)x^j\Bigr]\frac{1+tx}{1-x}.
\]
Extracting $[x^n]$: $\sum_{j<n}F(j)(1-(-t)^j)(1+t)+F(n)(1-(-t)^n)$;
dividing by $1+t$ and extracting $[t^k]$ gives, for $1\le k\le n-1$,
$-(-1)^kF(k)+(-1)^kF(n)$, and $\sum_{j\le n}F(j)$ for $k=0$.  Hence,
by the Fourier inversion of Theorem~\ref{thm:exactonepoint},
$\E[G_F(\sigma_m)]=\sum_{j\le n}F(j)+\sum_{k=1}^{n-1}r_k^{\,m}(F(n)-F(k))
=\sum_{j<n}F(j)(1-r_j^{\,m})+F(n)\sum_{k=0}^{n-1}r_k^{\,m}$.  Since
$G_F(\sigma)/n=\E_U[F(|C_\sigma(U)|)]$ for a uniform point $U$, this is
the finite-$n$ law; $F(j)=j^{r-1}$ gives the hook inner products, and
$j=nx$, $r_j^{\,m}\to(1-2x)^m$ gives the continuum statements
(Lemma~\ref{lem:splitmerge} for the identification).
\end{proof}

\emph{Status: PROVED.  VERIFIED: the hook identity
$(-1)^k\langle\chi_k,\sum c^r\rangle=n^{r-1}-k^{r-1}$ by brute-force
character sums for $r\le6$, $n\le9$ (\texttt{s10\_hooks34.py}); the
continuum law and the moments $r=2,3,4$ against simulation
(\texttt{s10\_sizebiased.py}, $2\times10^5$ trials): at $m=8$ the ten
decile masses agree to $\pm0.001$, the atom $0.1115$ vs $1/9$, and
$\E\sum p^{2,3,4}=0.5561,0.3947,0.3140$ vs $0.5556,0.3939,0.3131$.
Remarks.  (i) $r=2$ is Theorem~\ref{thm:exactonepoint}; the atom is
$\Pr[B_m=1]$ of Theorem~\ref{thm:blockcount}, as it must be.  (ii) The
drift identity $\E[\Delta L]=2L^2-\tfrac73\sum p_i^4$ ($L=\sum p_i^2$)
is now closed in expectation along the whole approach to
$\mathrm{PD}(1)$: $\E\sum p_i^4(m)=\tfrac14+\int_0^1(1-x^3)(1-2x)^mdx$.
(iii) Small marked blocks: $\Pr[\text{mass of }u\text{'s curve}\le\eta]
=\int_0^\eta(1-(1-2x)^m)dx\le\min(\eta,\,m\eta^2)$ --- the
$m\eta^2$ regime is the creation flux of Lemma~\ref{lem:recharge}(iii)
made exact for the block through a fixed point, and the crossover at
$\eta\asymp1/m$ is the trap scale.  (iv) The density $1-(1-2x)^m$ is
$\ge1-e^{-2mx}$ on $(0,\tfrac12]$: the size-biased pick is macroscopic with probability
bounded below uniformly in $m\ge1$ --- the quantitative handle for
recharge that TODO~10a(A)(iii) asked for, now superseded by
Corollary~\ref{cor:horizon}.}

\subsubsection{Recharge under the block bound}\label{sec:recharge}

Fix $K\ge3$ and let $H=H_T$ be the event of Corollary~\ref{cor:horizon}
with $K_\gamma(T)$ replaced by $K$.

\begin{lemma}[recharge of the marked total]\label{lem:srecharge}
Let $K\ge3$, $s_0\ge\varepsilon$ with $\varepsilon\le1/(2K)$.  Then for every
$T\ge1$,
\[
\Pr\Bigl[\max_{t\le T}s_t<\tfrac1{2K}\Bigr]
\;\le\;e^{-\varepsilon T/K}+2\varepsilon^2T+\Pr[H^c].
\]
The same holds for each marked coordinate separately: if
$a_0\ge\varepsilon$ then
$\Pr[\max_{t\le T}a_t<1/(2K)]\le e^{-\varepsilon T/(3K)}+2\varepsilon^2T+\Pr[H^c]$,
and likewise for $b$.
\end{lemma}

\begin{proof}
On $H\cap\{s_t<1/(2K)\}$ the environment has mass $>1-1/(2K)\ge\tfrac12$
in fewer than $K$ blocks, so it has a block of mass $z\ge1/(2K)$; the
chord with one cut in the marked material and one in that block has
probability $2s_tz\ge\varepsilon/K$ as long as $s_t\ge\varepsilon$, and
it raises $s$ by $z\ge1/(2K)$.  Flips and merges never decrease $s$, and
$s$ enters $(0,\varepsilon]$ from above with probability $\le2\varepsilon^2$
per step (Lemma~\ref{lem:strap}); so the event $\{\max_{t\le T}s_t<1/(2K)\}$
is contained in $\{s_t<\varepsilon\text{ for some }t\le T\}\cup H^c\cup
\{T\text{ consecutive failures of an event of probability }\ge\varepsilon/K\}$.
For a single coordinate $a$: if $s_t\le\tfrac12$ the same merge
argument applies to $a$ (probability $\ge2a\cdot1/(2K)\ge\varepsilon/K$);
if $s_t>\tfrac12$ and $a_t<1/(2K)$ then $b_t\ge\tfrac13$, a flip has
probability $2a_tb_t\ge2\varepsilon/3$ and lands at
$a'=\alpha+\beta\ge\beta\sim U(0,b_t)$, so $a'\ge1/(2K)$ with conditional
probability $\ge1-3/(2K)\ge\tfrac12$: success probability
$\ge\varepsilon/3\ge\varepsilon/(3K)$ per step in both regimes; the
entry flux into $(0,\varepsilon]$ is $\le2\varepsilon^2$ per step
(shrink of $a$, flip landing small; Lemma~\ref{lem:recharge}(iii)).
\end{proof}

\emph{Status: PROVED.  With $T=(3K/\varepsilon)\log(1/\delta)$ the
failure probability is $\delta+6\varepsilon K\log(1/\delta)+\Pr[H^c]$,
small precisely when $\varepsilon\ll1/K$ --- the regime in which
recharge is needed at all (for $\varepsilon\ge1/(2K)$ the start is
already log-macroscopic).  At the optimising
$\varepsilon=d^{-(1+2\alpha)/(2+2\alpha)}\in[d^{-3/4},d^{-1/2})$ of
Corollary~\ref{cor:kappafree} ($\alpha<1$) and $K=O(\log d)$,
$T=O(\varepsilon^{-1}\log^2d)\ll d$ and the failure probability is
$O(d^{-\gamma})+O(\varepsilon\log^2d)$, which is
$o\bigl((\varepsilon d)^{-2\alpha}\bigr)=o(d^{-\alpha/(1+\alpha)})$ since
$\varepsilon^{1+2\alpha}d^{2\alpha}=d^{-1/(2+2\alpha)}$.}

\begin{lemma}[log-drift of the marked total]\label{lem:logdrift}
On $H$, for a state with $s\le\tfrac12$,
\[
\E[\log s_{t+1}-\log s_t\mid\Gamma_t]\;\ge\;\frac{s}{K}\log\Bigl(1+\frac1{2Ks}\Bigr)-\frac{s^2}2 ,
\]
which is $\ge s\,(\log2-\tfrac14)/K>0$ for $s\le1/(2K)$; the negative
jumps of $\log s$ are those of $\tfrac12\log U$ (exponential tails),
the positive ones are unbounded above only through merges.
\end{lemma}

\begin{proof}
Merges: with probability $\ge2sz$ the total gains $z$, and the largest
environment block has $z\ge(1-s)/K\ge1/(2K)$; $\log(1+z/s)$ is
increasing in $z$.  Shrinks: $a\to a\sqrt U$ with probability $a^2$
changes $\log s$ by $\log(1-(a/s)(1-\sqrt U))\ge\log\sqrt U$, of mean
$-\tfrac12$; likewise $b$; $a^2+b^2\le s^2$.  Flips preserve $s$;
environment-only moves do not touch it.
\end{proof}

\emph{Status: PROVED (a calculation).  It says that, once the block
count is bounded, the marked total is a one-dimensional process with
positive log-drift below the level $\asymp1/K$ and light-tailed
down-jumps; standard drift arguments (Lyapunov function
$s^{-\theta}$, small $\theta$) then bound the occupation time of
$\{s<c/K\}$ over any window by a small fraction with polynomially
small failure probability.  This is the ``marked pair health''
input for the tail of the first-flip time in
\S\ref{sec:markovrenewal}; it is NOT written out.}

\subsubsection{The Markov-renewal form of Problem~\ref{prob:survival}}
\label{sec:markovrenewal}

Let $\tau$ be the first flip time of the quotient chain and, for
$|z|\le1$ and bounded measurable $F$,
\[
(K_zF)(x):=\E_x\bigl[z^{\tau}F(\Gamma_\tau)\bigr],
\qquad \|K_z\|_{L^\infty\to L^\infty}\le\sup_x\E_x|z|^{\tau}\le|z| ,
\]
since $\tau\ge1$.  $K_1$ is the post-flip kernel (Markov whenever
$\tau<\infty$ a.s.).

\begin{proposition}[generating function of the parity]\label{prop:genfun}
For $|z|<1$, as bounded functions of the initial state,
\[
\sum_{d\ge0}\Phi_d\,z^d\;=\;\frac1{1-z}\,(I+K_z)^{-1}(I-K_z)\mathbf 1 .
\]
\end{proposition}

\begin{proof}
With flip times $\tau_1<\tau_2<\cdots$, $\Phi_d=\sum_{k\ge0}(-1)^k
\Pr_x[\tau_k\le d<\tau_{k+1}]$ and
$\sum_dz^d\mathbf 1\{\tau_k\le d<\tau_{k+1}\}=(z^{\tau_k}-z^{\tau_{k+1}})/(1-z)$;
by the strong Markov property at the flip times
$\E_x[z^{\tau_k}]=(K_z^k\mathbf 1)(x)$.  Sum the Neumann series, which
converges since $\|K_z\|\le|z|<1$.
\end{proof}

\emph{Status: PROVED.  It is exact and separates the problem: the
factor $(I-K_z)\mathbf 1=\E_\cdot[1-z^{\tau}]$ carries the first-flip
tail (traps), the resolvent $(I+K_z)^{-1}$ carries the cancellation
(the parity of the number of flips per unit time).  Note
$(I-K_1)\mathbf 1=0$ wherever $\tau<\infty$ a.s.\ (on $\{ab>0\}$, since
$\sum_t2a_tb_t=\infty$ a.s.\ there --- taken for granted below), so the
right side is a difference quotient at $z=1$: the decay of $\Phi_d$ is the regularity of
$G(z):=(I+K_z)^{-1}(I-K_z)\mathbf 1$ on the closed disc.}

\begin{proposition}[conditional decay]\label{prop:condecay}
Let $\mathcal B$ be a Banach space of functions on the state space
containing $\mathbf 1$, on which the evaluations $F\mapsto F(x)$,
$x\in S'$, are uniformly bounded, and let $\gamma\in(0,1)$.  Suppose
\begin{enumerate}[leftmargin=2em]
\item[(R1)] $z\mapsto K_z\in\mathcal L(\mathcal B)$ is $C^{1,\gamma}$
on $\{|z|\le1\}$ (for a weighted sup-norm space
$\mathcal B=\{F:\sup|F|/V<\infty\}$ it suffices that
$\sup_x\E_x[\tau^{1+\gamma}V(\Gamma_\tau)]/V(x)<\infty$, by
$|\tau z_1^{\tau-1}-\tau z_2^{\tau-1}|\le2\tau^{1+\gamma}|z_1-z_2|^\gamma$);
\item[(R2)] $I+K_z$ is invertible on $\mathcal B$ for every $|z|\le1$,
with $\sup_{|z|\le1}\|(I+K_z)^{-1}\|<\infty$.
\end{enumerate}
Then $\sup_{x\in S'}|\Phi_d(x)|\le C\,d^{-\gamma}$.
\end{proposition}

\begin{proof}
Under (R1)--(R2), $G$ is $C^{1,\gamma}$ on the closed disc with
$G(1)=0$, so $\hat\Phi(z)=(G(z)-G(1))/(1-z)$ extends to a
$C^{0,\gamma}$ function on the closed disc; it is analytic inside, so
its boundary Fourier coefficients are the $\Phi_d$, and the Fourier
coefficients of a H\"older-$\gamma$ function are $O(d^{-\gamma})$.
Evaluate at $x\in S'$.
\end{proof}

\emph{\textbf{SESSION-11 CORRECTION.}  Condition (R2) is FALSE at
$z=\pm1$ for every function space containing the bounded functions:
$K_1\varepsilon=-\varepsilon$, because exactly one flip has occurred at
time $\tau$ (Proposition~\ref{prop:r2false}).  Proposition~\ref{prop:condecay}
is therefore vacuous as stated.  The same holds in the
label-free bookkeeping, with $\kappa_{-1}[(-1)^E]=-(-1)^E$.  It is
repaired in \S\ref{sec:s11}: (R2) is replaced by the two conditions
(C1) and (C2) of \S\ref{sec:s11left}, both statements about the single
Markov kernel $\kappa_1$, and everything on the closed disc away from
$z=\pm1$ is proved unconditionally by
Theorem~\ref{thm:idleloop}.  The discussion of (R1) below is
unaffected, and so is the reduction paragraph at the end of this
subsection.  Read the rest of this paragraph with that substitution.}

\emph{Status: PROVED as a reduction; (R1) and (R2) are OPEN.  What each
needs, and what is in hand:}

\emph{(R1) --- traps.  A uniform bound $\sup_{x\in S'}\Pr_x[\tau>t]\le
C_Kt^{-2}$ (measured: $t^2\Pr[\tau>t]\to11$ from stationary
environments, \S\ref{sec:parity}(c)) gives every moment $\tau^{1+\gamma}$,
$\gamma<1$, hence $\gamma$ arbitrarily close to $1$ in
Proposition~\ref{prop:condecay}, far above the $4/9$ that
Corollary~\ref{cor:kappafree} needs.  Ingredients: the state-independent
trap-creation flux (Lemma~\ref{lem:recharge}(iii)); the recharge of
Lemma~\ref{lem:srecharge} and the drift of Lemma~\ref{lem:logdrift},
both available on $H$ by Corollary~\ref{cor:horizon}; and the fact
that a trap of depth $g$ is left by a flip alone, at rate
$2g(s-g)$, needing no environment.  The proof of the $1/t^2$ tail is
a renewal computation with these inputs (the subcritical kernel
$\approx2/(s^2m^2)$ of \S\ref{sec:parity}(c)), which we have not
written.}

\emph{(R2) --- cancellation.  This is the parity content and cannot be
had from trap accounting.  It is needed on the whole closed disc: the
trivial bound $\|K_z\|\le|z|$ gives nothing on $|z|=1$.  At $z=1$ it
is the statement that $-1$ is not in the spectrum of the post-flip
chain $K_1$ on $\mathcal B$, with a uniform inverse; near $z=1$ it
follows by perturbation under (R1).  At other points of the circle it
is the non-lattice condition of Markov renewal theory --- no $F\ne0$
with $\E_x[z^{\tau}F(\Gamma_\tau)]=-F(x)$ --- and at $z=-1$, where
$K_{-1}F=\E_\cdot[(-1)^{\tau}F(\Gamma_\tau)]$, it says the parity of the
inter-flip times is not asymptotically deterministic; this is where
the $(-1)^d$ alternation of \S\ref{sec:parity}(a) lives and it needs
its own statement.  The Doeblin structure is on the marked pair: two consecutive
flips from any $(a,b)$ with total $s$ produce $a''$ whose law dominates
$\tfrac14\cdot\mathrm{Unif}[s/4,3s/4]$ (the trapezoid of
Lemma~\ref{lem:recharge}(ii) is symmetric and unimodal, so the first
flip lands in $[s/4,3s/4]$ with probability $\ge\tfrac12$, and from
there the second has density $\ge1/s$ on that interval).  This is a
minorisation on the fibre $\{a'+b'=s,\ \mathrm{env}\}$ only; the base
$(s,\mathrm{env})$ is not regenerated, so $\mathcal B$ must be a
weighted space in which the post-flip chain is $V$-geometrically
ergodic in the base coordinates as well (Meyn--Tweedie), with the
weight matched to the moment condition in (R1).  On $H$ the base is
finite-dimensional ($<K$ blocks), which is what makes this plausible.
Because the label parity is determined by the environment block-count
parity (Lemma~\ref{lem:parity}), no argument can avoid the environment
in (R2) either; but the environment enters only through the ergodicity
of the base chain between flips, not through any fine property of
$\mathrm{PD}(1)$.}

\emph{Reduction to Problem~\ref{prob:survival}.  Take
$S'=S'_K:=\{a,b\ge1/(16K)\}$ and enforce $H_{9d}$ (so $C=C(K)$ in
Proposition~\ref{prop:condecay} may depend on the block bound, which
holds up to time $d$).  Granting Proposition~\ref{prop:condecay} on
$S'_K$ with $C(K)$ polynomial in $K$, and granting the occupation-time
consequence of Lemma~\ref{lem:logdrift} (that after $s$ first reaches
$1/(2K)$ it stays $\ge1/(4K)$ for at least half of the next $T_2$
steps, except with polynomially small probability; NOT written out),
the chain from $\Gamma_0\sim\mu_j$ with $a_0,b_0\ge\varepsilon$ enters
$S'_K$ at a stopping time $T'$: first $s$ reaches $1/(2K)$ within
$T_1=O(\varepsilon^{-1}K\log d)$ steps (Lemma~\ref{lem:srecharge});
then, on the times with $s_t\ge1/(4K)$ and $\min(a_t,b_t)\ge\varepsilon$
(the latter failing with probability $\le4\varepsilon^2(T_1+T_2)$ by
the flux bound), a flip occurs with probability $\ge2\varepsilon(s_t-\varepsilon)
\ge\varepsilon/(4K)$ per step and lands with $\min(a',b')\ge s_t/4\ge1/(16K)$
with conditional probability $\ge\tfrac12$ (the trapezoid is
symmetric and unimodal), so $T'\le T_1+T_2$ with
$T_2=O(\varepsilon^{-1}K\log d)$ except with probability
$O(d^{-\gamma})+O(\varepsilon K\log d)$.  Note that a single
application of Lemma~\ref{lem:srecharge} to each coordinate does not
suffice: it makes $a$ and $b$ large at different times, and a
coordinate of size $1/(2K)$ shrinks on the scale $K^2\ll\varepsilon^{-1}K$.
By the strong Markov property
$|\Phi_d(\Gamma_0)|\le\E|\Phi_{d-T'}(\Gamma_{T'})|+\Pr[T'>T_1+T_2]
\le C(K)(d/2)^{-\gamma}+O(d^{-\gamma})+O(\varepsilon K\log d)$, and
with $K=O(\log d)$ and $\varepsilon$ as in Corollary~\ref{cor:kappafree}
this is $\le C(\varepsilon d)^{-2\alpha}$ for $\alpha<\gamma/2$ (for
$\varepsilon\ge1/(16K)$ no recharge is needed).  Thus
\emph{Problem~\ref{prob:survival}, hence the iid model, is reduced to
(R1) and (R2) on the log-macroscopic set $S'_K$ plus the
occupation-time step}, with the environment controlled throughout by
Corollary~\ref{cor:horizon}.}
